\documentclass{situationist-preprint}

\title{Astrology and the Absence of Error-Correction}
\msid{SCI-01-b}
\msfamily{SCI-01}
\msdiscipline{Philosophy of Science}
\mstype{Companion manuscript}
\mskeywords{astrology; demarcation; pseudoscience; error-correction; falsifiability}

\begin{document}
\maketitle

\begin{abstract}
SCI-01-a argues that science is the enterprise with institutionalized error-correction. Astrology is the standard test case for any criterion of demarcation, and it has embarrassed the classical ones: its claims are testable and have been tested, so falsifiability classifies it with science. We argue that the present criterion classifies it correctly and for the right reason. The checks of astrology's central principles took place in another enterprise, and the astrological canon survived them unaltered; within astrology the refutation schedule of those principles is empty. The revisions astrologers do make, such as the assignment of rulerships to newly discovered planets, add to the periphery without removing anything from the core.
\end{abstract}

\section{The case}\label{s:case}

Astrology holds that the positions of the sun, moon and planets at a person's birth bear on that person's character and fortunes, and it codifies the bearing in a doctrine of signs, houses and aspects largely fixed by the second century. \citet{Popper1963} placed astrology on the wrong side of his line on the ground that its practitioners evade refutation by vague prophecy. The ground is insufficient. Many astrological claims are precise enough to be tested, and they have been. In the double-blind test of \citet{Carlson1985}, astrologers who had committed in advance to the procedure were unable to match natal charts to personality profiles better than chance. By the falsifiability criterion, then, astrology is a falsified science, not a non-science: the verdict Laudan accepted for creation science, and the one that SCI-01-a rejects for both.

\citet{Thagard1978} located the defect instead in the absence of progress and in the indifference of the astrological community to anomalies and to rival theories. Thagard's diagnosis points in the right direction, but it rests on features of the community's attitudes, and attitudes vary from one practitioner to another. SCI-01-a locates the defect in the structure of the enterprise.

\begin{proposition}\label{p:astro}
Astrology is not a science, and its central principles have an empty refutation schedule within it.
\end{proposition}

\begin{proof}
By Definition~3.2 of SCI-01-a, the refutation schedule of an enterprise for a claim is the set of times at which its own checks would refute the claim. The tests of astrology, from Carlson's study to the statistical surveys that preceded it, were designed, conducted and published by physicists and psychologists in the enterprises of physics and psychology, where their results entered the record. No procedure internal to astrology took their results as input, and no principle of the canon was withdrawn in consequence. The schedule for those principles within astrology is therefore empty, and by Corollary~4.6 of SCI-01-a an enterprise whose checks of a false claim cease, here never having begun, retains it forever and is not a science.
\end{proof}

The proof shows why testability was the wrong test. The question is not whether a claim can be checked by someone, but whether the enterprise that holds it has standing procedures by which a failed check removes it from its record. Astrology's claims are checkable, and checked; astrology is not self-checking.

\section{Objections and replies}\label{s:objections}

\begin{objection}[Revision]
Astrology does change. When Uranus, Neptune and Pluto were discovered, astrologers assigned them rulerships and meanings; schools dispute house systems and the choice between the tropical and the sidereal zodiac.
\end{objection}

\begin{reply}
Every one of these changes adds to the doctrine or chooses among variants of it, and none removes a principle. The tropical zodiac is the clearest case. Once the precession of the equinoxes had detached the signs from the constellations that name them, the tradition did not withdraw the claim that the sun's sign at birth bears on character; it redefined the signs so that the claim survives whatever the sky does. Corollary~4.6 of SCI-01-a requires for science that no claim, central or peripheral, be exempt from removal, and here the central claims are exempt by construction.
\end{reply}

\begin{objection}[Ancient astronomy]
In antiquity astrology and astronomy were practiced by the same people, Ptolemy among them. If they were one enterprise, astrology was once part of a science.
\end{objection}

\begin{reply}
The criterion classifies enterprises at times, and it is open to the conclusion that the boundary has moved. But the question is whether the enterprise had standing checks that could remove the astrological principles, and it did not: the positions predicted by the astronomical tables were corrected against observation, while the doctrine connecting positions to human affairs was justified by analogy and authority and was never exposed to the same correction. The single practitioner housed two records, one self-cleansing and one not.
\end{reply}

\section{Conclusion}\label{s:conclusion}

Astrology is neither untestable nor untested. It is unscientific because the tests of its principles happened elsewhere, in enterprises that recorded the failures, and left its own canon unaltered. That is the absence of institutionalized error-correction which SCI-01-a identifies as the mark of non-science.

\begin{thebibliography}{9}
\bibitem[Carlson(1985)]{Carlson1985} Carlson, S. (1985). A double-blind test of astrology. \emph{Nature}, 318, 419--425.
\bibitem[Popper(1963)]{Popper1963} Popper, K.~R. (1963). \emph{Conjectures and Refutations: The Growth of Scientific Knowledge}. Routledge and Kegan Paul.
\bibitem[Thagard(1978)]{Thagard1978} Thagard, P.~R. (1978). Why astrology is a pseudoscience. In P.~D. Asquith and I.~Hacking (eds.), \emph{PSA 1978}, vol.~1, 223--234. Philosophy of Science Association.
\end{thebibliography}

\end{document}
